\subsection{ধারণাভিত্তিক MCQ}
প্রচলিত সংখ্যাগত শ্রেণি, হিসাব এবং সাধারণ বিভ্রান্তি মিলিয়ে তৈরি নিজস্ব অনুশীলন। উত্তরমালা শেষে আছে।
\begin{multicols}{2}\begin{enumerate}
\item\chTwoMCQBlock{প্রচলিত HSC শ্রেণিতে Narrowband-এর পরিসর কোনটি?
\par ক. 45–300 bps\par খ. 1200–9600 bps\par গ. 1–10 Mbps\par ঘ. 45–300 MB/s}
\item\chTwoMCQBlock{প্রচলিত Voiceband শ্রেণির পরিসর কোনটি?
\par ক. 45–300 bps\par খ. 1200–9600 bps\par গ. 1–100 Mbps\par ঘ. 1200–9600 Byte/s}
\item\chTwoMCQBlock{500 Byte দুই সেকেন্ডে পৌঁছালে গড় ডেটা হার কত?
\par ক. 250 bps\par খ. 4000 bps\par গ. 2000 bps\par ঘ. 8000 bps}
\item\chTwoMCQBlock{4 Mbps সংযোগে দশমিক এককে 10 MB ফাইলের আদর্শ প্রেরণসময় কত?
\par ক. 2.5 s\par খ. 10 s\par গ. 80 s\par ঘ. 20 s}
\item\chTwoMCQBlock{একই পথে 25, 8 ও 40 Mbps সংযোগ থাকলে কোনটি গতির সীমা নির্ধারণ করে?
\par ক. 8 Mbps\par খ. 25 Mbps\par গ. 40 Mbps\par ঘ. 73 Mbps}
\item\chTwoMCQBlock{পরপর প্যাকেটের পৌঁছানোর বিলম্বের ওঠানামাকে কী বলে?
\par ক. Throughput\par খ. Jitter\par গ. Bandwidth\par ঘ. Frequency}
\item\chTwoMCQBlock{8 ডেটা + 1 স্টার্ট + 1 স্টপ বিটের ফ্রেমের কার্যদক্ষতা কত?
\par ক. 72.73\%\par খ. 88.89\%\par গ. 80\%\par ঘ. 100\%}
\item\chTwoMCQBlock{8 ডেটা + 1 স্টার্ট + 1 প্যারিটি + 1 স্টপ বিটের ফ্রেমের কার্যদক্ষতা কত?
\par ক. 80\%\par খ. 90\%\par গ. 66.67\%\par ঘ. 72.73\%}
\item\chTwoMCQBlock{80 Byte ডেটা ও মোট 4 Byte বাড়তি নিয়ন্ত্রণতথ্যের ব্লকের কার্যদক্ষতা প্রায় কত?
\par ক. 95.24\%\par খ. 90\%\par গ. 80\%\par ঘ. 72.73\%}
\item\chTwoMCQBlock{1100 bps লাইনে প্রতি ফ্রেমে 8 ডেটা+1 স্টার্ট+2 স্টপ বিট থাকলে উপযোগী ডেটা হার কত?
\par ক. 1100 bps\par খ. 800 bps\par গ. 1000 bps\par ঘ. 880 bps}
\item\chTwoMCQBlock{প্রাপকের সঠিক সময়ে বিট শনাক্ত করার সমন্বয় কোনটি?
\par ক. Roaming\par খ. Routing\par গ. Bit Synchronization\par ঘ. Resource Pooling}
\item\chTwoMCQBlock{IP Multicast শনাক্ত করার প্রধান বৈশিষ্ট্য কোনটি?
\par ক. শুধু একজন প্রাপক\par খ. সকল গ্রাহককে আলাদা Email\par গ. সব Port-এর Signal\par ঘ. নিবন্ধিত Group Address}
\item\chTwoMCQBlock{Thinnet-এর ঐতিহাসিক Ethernet নাম কোনটি?
\par ক. 10BASE2\par খ. 10BASE5\par গ. 100BASE-TX\par ঘ. 1000BASE-T}
\item\chTwoMCQBlock{চার জোড়াবিশিষ্ট Ethernet UTP কেবলে কতটি পরিবাহী তার থাকে?
\par ক. 2\par খ. 8\par গ. 4\par ঘ. 16}
\item\chTwoMCQBlock{পরীক্ষায় প্রচলিত Single-mode Fiber-এর কোরের ব্যাসের পরিসর কোনটি?
\par ক. 50–100 $\mu$m\par খ. 185–500 $\mu$m\par গ. 8–12 $\mu$m\par ঘ. 8–12 mm}
\item\chTwoMCQBlock{ফাইবারে পূর্ণ অভ্যন্তরীণ প্রতিফলনের জন্য কোন সম্পর্ক প্রয়োজন?
\par ক. কোরের প্রতিসরণাঙ্ক $<$ ক্ল্যাডিংয়ের প্রতিসরণাঙ্ক\par খ. দুটি প্রতিসরণাঙ্ক অবশ্যই সমান\par গ. কোর অবশ্যই তামার তৈরি\par ঘ. কোরের প্রতিসরণাঙ্ক $>$ ক্ল্যাডিংয়ের প্রতিসরণাঙ্ক}
\item\chTwoMCQBlock{HSC-তে ব্যবহৃত GEO উপগ্রহের আনুমানিক উচ্চতা কত?
\par ক. 36,000 km\par খ. 360 km\par গ. 3600 km\par ঘ. 360,000 km}
\item\chTwoMCQBlock{ছোট স্যাটেলাইট ভূকেন্দ্রের সংক্ষিপ্ত নাম কোনটি?
\par ক. VLAN\par খ. VSAT\par গ. VoLTE\par ঘ. UTP}
\item\chTwoMCQBlock{Classic Bluetooth Piconet-এ কেন্দ্রসহ সর্বোচ্চ সক্রিয় যন্ত্র কতটি?
\par ক. 7\par খ. 255\par গ. 8\par ঘ. 256}
\item\chTwoMCQBlock{Wi-Fi/WLAN-এর মানের পরিবার কোনটি?
\par ক. IEEE 802.16\par খ. IEEE 802.3\par গ. IEEE 802.15.1\par ঘ. IEEE 802.11}
\item\chTwoMCQBlock{WMAN/WiMAX-এর মানের পরিবার কোনটি?
\par ক. IEEE 802.16\par খ. IEEE 802.11\par গ. IEEE 802.3\par ঘ. ASCII}
\item\chTwoMCQBlock{FOMA কোন প্রজন্মের সেবার ঐতিহাসিক নাম?
\par ক. 1G\par খ. 3G\par গ. 2G\par ঘ. 5G}
\item\chTwoMCQBlock{সচল কল এক বেস স্টেশন থেকে অন্যটিতে স্থানান্তরের প্রক্রিয়া কী?
\par ক. Roaming\par খ. Multicast\par গ. Handoff\par ঘ. Modulation}
\item\chTwoMCQBlock{একই ক্যাম্পাসের কয়েকটি ভবনের LAN যুক্ত হলে বিস্তৃতির কোন পরিচয় প্রযোজ্য?
\par ক. PAN\par খ. WAN\par গ. Peer-to-peer\par ঘ. CAN}
\item\chTwoMCQBlock{সুইচ প্রাপ্ত ফ্রেমের কোন ঠিকানা দেখে MAC Table শেখে?
\par ক. Source MAC\par খ. Destination IP\par গ. Source URL\par ঘ. Destination Domain}
\item\chTwoMCQBlock{গন্তব্যের MAC অজানা হলে সাধারণ সুইচ কী করতে পারে?
\par ক. সব ফ্রেম বাতিল\par খ. আগত পোর্ট ছাড়া সংশ্লিষ্ট অন্য ফরওয়ার্ডিং পোর্টে ছড়িয়ে দেওয়া\par গ. শুধু আগত পোর্টে ফেরত দেওয়া\par ঘ. প্রোটোকল রূপান্তর}
\item\chTwoMCQBlock{স্থানীয় নেটওয়ার্কের বাইরের গন্তব্যে পাঠাতে কম্পিউটারের Default Gateway সাধারণত কী?
\par ক. অ্যানালগ বর্ধক\par খ. শুধু প্রোটোকল রূপান্তরকারী\par গ. বাইরের নেটওয়ার্কে যাওয়ার রাউটার\par ঘ. শুধু রিপিটার}
\item\chTwoMCQBlock{6-নোডের Full Mesh-এ দ্বিমুখী লিংক কতটি?
\par ক. 6\par খ. 12\par গ. 30\par ঘ. 15}
\item\chTwoMCQBlock{6-নোডের Full Mesh-এ প্রতি নোডের পোর্ট কতটি?
\par ক. 5\par খ. 6\par গ. 15\par ঘ. 30}
\item\chTwoMCQBlock{Full Mesh-এ 45টি লিংক থাকলে নোড কতটি?
\par ক. 9\par খ. 10\par গ. 15\par ঘ. 45}
\item\chTwoMCQBlock{ভার্চুয়াল সার্ভারে নিজে OS ও অ্যাপ্লিকেশন চালানোর সেবা কোন মডেলের?
\par ক. SaaS\par খ. PaaS\par গ. IaaS\par ঘ. Public}
\item\chTwoMCQBlock{ল্যাপটপে ফাইল মুছে ফেলায় ক্লাউডের সিঙ্ক্রোনাইজড কপিও মুছে গেছে। কোন সিদ্ধান্তটি সঠিক?
\par ক. সিঙ্ক্রোনাইজেশন সবসময় ব্যাকআপ\par খ. ক্লাউডে ফাইল মুছতে পারে না\par গ. Private Cloud হলে পাসওয়ার্ড লাগে না\par ঘ. আলাদা ব্যাকআপ বা আগের সংস্করণ প্রয়োজন}
\end{enumerate}\end{multicols}
\subsubsection*{উত্তরমালা}
\begin{multicols}{2}
১. ক, ২. খ, ৩. গ, ৪. ঘ, ৫. ক, ৬. খ, ৭. গ, ৮. ঘ।\par
৯. ক, ১০. খ, ১১. গ, ১২. ঘ, ১৩. ক, ১৪. খ, ১৫. গ, ১৬. ঘ।\par
১৭. ক, ১৮. খ, ১৯. গ, ২০. ঘ, ২১. ক, ২২. খ, ২৩. গ, ২৪. ঘ।\par
২৫. ক, ২৬. খ, ২৭. গ, ২৮. ঘ, ২৯. ক, ৩০. খ, ৩১. গ, ৩২. ঘ।\par
\end{multicols}
